\textbf{COLUNI 2023 -- Questão 34}\par\medskip Leia o texto a seguir:

\par

“Quando, a bordo do Santos, chegaram Carlos e Alfredo a Paranaguá, encantou-os a vista da terra. Abria-se aos seus olhos um lindo panorama: o porto, vasto e quieto, cercado de um amplo círculo de montanhas. Alfredo estava ansioso por descer à terra: − Vamos, vamos! Quero ficar conhecendo a capital de mais um Estado do Brasil! − Que? − atalhou o irmão sorrindo − então Paranaguá é capital? − É verdade! − exclamou o pequeno − que tolice a minha! [...] Baixaram à terra, e com eles um homem alto, forte, louro, que se ofereceu para lhes mostrar a cidade. Era um alemão, que, apesar de não estar no Brasil há mais de cinco anos, já falava perfeitamente o português. Chamava-se Schumann, e era muito conversador e simpático. Os quatro viajantes percorreram com prazer a pequena cidade, conversando. O alemão residia na capital de Santa Catarina, mas conhecia bem o Paraná. − Não podem fazer ideia do que é a estrada de ferro que vai daqui até a Lapa e Rio Negro. É uma maravilha que honra a engenharia brasileira. Tem obras de arte admiráveis, viadutos, túneis. Em certos pontos, na serra, o trem fica suspenso sobre abismos, cuja contemplação dá vertigens. [...] − E o Estado é muito rico? − perguntou Carlos. − Muito rico, não é, − disse Schumann − mas é próspero, e as bases da sua prosperidade são os pinhais e o mate. O pinheiro do Paraná, araucária brasiliensis, dá uma madeira tão boa como a da Noruega: é uma árvore corpulenta, que às vezes atinge à altura de trinta e seis metros. Há aqui pinheirais admiráveis, em grandes florestas que cobrem os planaltos. Li há poucos dias, um artigo, cujo autor calcula em mais de oitenta milhões os pinheiros que formam essas florestas... [...] Assim conversando e passeando, passaram os viajantes cerca de hora e meia em Paranaguá. Voltaram para bordo. O Santos tomou de novo o rumo do sul.

\par

Fonte: BILAC, Olavo; BONFIM, Manuel. Através do Brasil. Francisco Alves, RJ, 1948, 36ª ed. Disponível em: https://www.unicamp.br/iel/memoria/Ensaios/LiteraturaInfantil/72atraves.htm. Acesso em: 26 ago. 2022.

\par

Considerando o texto acima e os conhecimentos sobre as regiões brasileiras, é INCORRETO afirmar que a região sul\begin{enumerate}
\item Ocupa áreas situadas nas menores latitudes brasileiras, por se localizarem abaixo do Trópico de
       Capricórnio.
\item É a menor região geográfica do Brasil em relação ao número de estados, cujas capitais são: Curitiba,
       Florianópolis e Porto Alegre.
\item Recebeu grande afluxo de migrantes europeus que marcaram as feições regionais como o padrão de uso
       da terra, as características arquitetônicas e os elementos culturais.
\item Apresenta elementos físicos bastante distintos das demais regiões, como o tipo climático, a cobertura
       vegetal e as feições do relevo.
\end{enumerate}
